% Noether P09 Korean producer draft — UNCHECKED
% T06-U51; ED0002 lines6946--6970; source 1098 LF B / D33127E20D06527574B7C11724CE13A84CB6506D4144F449DE878591828CE187
% Pointer v009 / ED0002: B06BE3530D9CF2E82B56FDBA7FE41D5D044DF2425DFA2A059D4939EAA2F7A6C2 / C9A125167ACB33D914EE4374B65AE7CDF0052F568371B8B77B720EA178ABF0E3
% Translation only; proof and degree terminology remain checker-held.

이제 영역 $\mathfrak G_\nu$들이 $[H]$의 다항식 $f(\eta)$ 이외에는 어떠한 공통 원소도 갖지 않음을 보이겠다. $z$를 $\eta$의 임의의 정수적 대수함수이되 유리함수는 아닌 것으로 하고, $[H]$에 대하여 기약인 다음 방정식을 만족한다고 하자.
\[
\begin{array}{@{}l@{\quad}c@{}}
(2)& z^\chi+a_1(\eta)z^{\chi-1}+a_2(\eta)z^{\chi-2}+\cdots+a_\chi(\eta)=0.
\end{array}
\]
따라서 $\chi>1$이다. 다항식 $a_1(\eta),\ldots,a_\chi(\eta)$의 차수 가운데 가장 큰 것을 $\lambda$라 하자. $z$가 $\mathfrak G_\nu$에 속하면 양 $\zeta=z-f(\eta)$도 $[H]$에 대하여 기약인 다음 방정식을 만족한다.
\[
\begin{array}{@{}l@{\quad}c@{}}
(3)& \zeta^\chi+b_1(\eta)\zeta^{\chi-1}+b_2(\eta)\zeta^{\chi-2}+\cdots+b_\chi(\eta)=0.
\end{array}
\]
$b_1(\eta),b_2(\eta),\ldots,b_\chi(\eta)$는 $\zeta$의 대칭함수이므로, (1)에 따라 그 최저차항의 차수는 각각 적어도 $\nu,2\nu,\ldots,\chi\nu$이다. (2)와 (3)을 비교하면
\[
 b_1(\eta)=a_1(\eta)+\chi f(\eta)
\]
를 얻는다.
